\documentclass[12pt,letterpaper]{article}
\usepackage{graphicx} % Required for inserting images
\usepackage{xparse}
\usepackage{xcolor}
\usepackage{bbm}
\usepackage{hyperref}
\usepackage{url}
\usepackage{amsmath,mathtools}
\usepackage{amsthm}
\usepackage[margin=1in]{geometry}
\usepackage{tikz}
\usetikzlibrary{matrix,positioning,calc}
\title{Response to reviewers}
\author{Victor Funes-Leal \and Jared Hutchins}
\date{}

\begin{document}
\maketitle

\section*{Reviewer 1}
\subsection*{Comment 1}
{\it
\textbf{Inbreeding as an externality?} The paper argues that the costs of inbreeding are an externality borne by the dairy industry as a whole, but that individual dairy bull breeders do not have an incentive to consider these costs in their individual
breeding decisions. The second section of the paper attempts to provide a
descriptive conceptual framework for this argument, repeatedly comparing
inbreeding to pesticide resistance management (specifically, Bt crops and refugia).
However, I find the comparison to Bt resistance tenuous and more broadly am not
yet convinced that the costs of inbreeding are an externality at all. Rather than using a descriptive analogy to argue their point, I think AJAE readers would be more convinced by a clear, concise theoretical model: Conceptually, pesticide resistance has been well-established through theoretical modeling (e.g. Hueth and Regev 1974; Miranowski and Carlson 1986; Cho et al. 2024) to present a type of common pool resource (CPR) problem for farmers, whereby any farmer may use pesticides
and draw upon the genetic pool of susceptibility in the pest population (nonexcludability), but in doing so decreases the pool available to all farmers using it in the future due to resistance selection (rivalry). These features are borne out and clearly illustrated in bioeconomic models of pesticide resistance. With dairy cattle inbreeding, the nature and existence of the externality is less clear: Where is nonexcludability (or nonrivalry) present? Based on the information in the paper, high discount rates could provide a simple alternative explanation for increased inbreeding: Breeders may simply favor the short-run trait gains over the long-run risks of inbreeding, even if these discounted future risks are internalized in the
market. Of course, dairy industry discount rates may likely be higher than is socially optimal (Wuepper et al. 2023), but this does not mean there is necessarily an externality present. One might also construe inbreeding as a type of pecuniary externality rather than a ``real'' externality, in which competition and market forces among dairy bull breeders have negative consequences—albeit internalized and market-transmitted—on dairy producers (limited variety of non-inbred bull semen, e.g. in the event of a negative homozygous trait coming to light).

Part of the difficulty I have seeing the externality argument lies in the vague
description of the risks of inbreeding and their costs. The main long-run risk the
authors describe in section 2 is an increased ``likelihood of unforeseen genetic traits that could harm the industry''... such as? That is all readers get, and they might be left wondering (as I was) why there’s not an incentive for improvements in GT technologies to better identify and predict these negative traits? A quick literature search (I have no background in dairy science or economics) revealed that such efforts may in fact be underway (Howard et al. 2017). This vagueness about the costs of inbreeding persists in Section 6 (Cost-Benefit Analysis), which uses an old estimate of the marginal costs of inbreeding (24 USD per 1\% increase in inbreeding), without any explanation of the basis for that figure and how it was calculated (does it depend on a discount rate?). Since it plays a pivotal role in the paper’s economic conclusions, this figure deserves more attention and validation (ideally, this paper would generate an updated estimate of the marginal cost of inbreeding, though I understand that may be beyond the paper’s scope).}

\subsection*{Comment 2}
{\it
\textbf{Empirical methods and identification}. Overall, the evidence provided seems pretty clear that GT dramatically increased inbreeding rates. However, I do have some concerns, questions, and suggestions regarding the econometric methods used. The paper adopts a frequently used DID-PSM approach, which is compelling, but their setup makes a subtle argument in the definition of the treatment group, with implications for exactly what is being causally estimated: The authors define
the treatment group of bulls as those coming from lineages of “superstar” founder
bulls from 1998-2004, those whose descendants represented a very high proportion
of all future bulls. They compare these superstar founders to those from matched
controls who were not superstars but had comparable (measured) genetic traits.
They then examine (using DID) how inbreeding rates changed between these two groups after the advent of GT technology. Notably, the treatment group is not defined as those bull lineages in which GT was adopted (the authors do not observe
this information, and argue that GT was so widespread as to make estimation of GT
treatment effects infeasible). So, what do these causal estimates tell us? Assuming
the PSM is successful, this makes selection of trait-matched lineages into superstar status as good as random. Within the subsample of superstar and matched
lineages, the treatment here must technically be defined as: GT (post-2010) x
(randomly selection into) superstar status. But to the extent that it also increased inbreeding among non-superstars, this estimation can’t reveal the aggregate effect of GT on inbreeding. Indeed, Figure 7 appears to show that inbreeding in the matched controls increased more rapidly post-GT (especially starting around 2014).

These subtle differences matter for the paper’s conclusions and economic
implications. Yet the BCA in section 6 seems to sweep these subtleties aside- the
superstar-specific GT effect is used to calculate the aggregate, cumulative effect of GT on inbreeding, which is not correct given the definition of the treatment in the econometric model. In particular, it omits the costs of increased inbreeding in nonsuperstars.}

\subsection*{Comment 3}
{\it
Trait improvement equation (1): I don’t think equation (1) serves any substantive
purpose in the paper. Economists readers won’t be able to make much of its utility (I wasn’t at least, and I have some knowledge of genetics, though none of dairy bull breeding)—the units in this equation are unclear to me, and it’s not linked to any theoretical model that might shed light on the incentives in the industry in relation to the alleged externality. Nor is it used in any of the econometric estimation (as far as I can tell).}

\subsection*{Comment 4}
{\it
Examining Figures 7 and 8, I wonder if we are observing that GT led to earlier
increases in inbreeding in superstars v. matched controls (perhaps due to earlier GT adoption?) with a possible leveling-out starting in 2019 (last year observed in data).

One could speculate that after 2019, the non-superstars’ inbreeding rates eventually caught up with those of superstars, perhaps as GT use further penetrated
down into more marginal lineages. Indeed, one might expect that breeders of the
control bulls—having similar genetic traits but less popular for purely random
reasons—might eventually realize they could use GT to more rapidly improve their
lines, communicate that value to dairy producers, and gain more market share.
}

\subsection*{Comment 5}
{\it Line-level clustered standard errors. First, please report the number of lines
(clusters) in the estimation table (table notes are also needed throughout- for *’s, standard errors, reference to corresponding model equations, etc.). Also, one could argue that firm-clustering may be desirable as well: You could do both (as long as lines weren’t fully contained within firms, or vice versa), using multiway clustering!}

\textbf{Reply}: Table 2 has been corrected; it includes the number of clusters included in the estimation of the cluster-robust covariance matrices for all four models. We have also included a new table in the appendix that provides for two-way clustered standard errors by sire ID and company ID.

\subsection*{Comment 6}
{\it
BCA in Section 6- discounting? It isn’t clear how the inbreeding USD costs and
genetic trait benefits are aggregated over the time period considered, and if a
discount rate was used in that calculation. Moreover, if the costs of inbreeding
include longer-term risks (see above comment), it isn’t clear how the discount rate
might affect these calculations. Your BCA warrants its own appendix section
detailing these calculations and assumptions for readers.}

\subsection*{Comment 7}
{\it PSM balance assessment \& alternative matching estimators. From my
understanding of the matching literature, raw t-stats and p-values in your balance
table are not the best way to assess balance; ‘mean standardized differences’
(MSD) seem to have gained prominence as a way of assessing balance, and I
personally much prefer a plot of MSDs, rather than tabular format (see Greifer
2025). Why did you only consider PSM in your matching strategy rather than
alternatives such as nearest-neighbor, or inverse-probability weighting on the
propensity score? If they haven’t already done so, I encourage the authors to read through the matching chapter of the Causal Inference Mixtape textbook
(Cunningham 2021). Also, considering that superstar status in founder bulls might
be related to firm-level marketing effort, can you include a firm-fixed effect in the propensity-score model? Even better, in my opinion, is a within-cluster(firm)
matching estimator, where matches are drawn from the same firm, if available,
before drawing matches from other firms (Arpino and Cannas 2016).}

\subsection*{Comment 8}
{\it In the main text’s description of the estimation models, some of the notation was hard to follow with incomplete definitions (e.g. $inb$, $yob$). I had to hunt around and deduce indirectly from the text to deduce that $k$ indexes firm and $\alpha_{k}$ was the firm-level fixed effect (there’s also a typo in the appendix, with $\alpha_{j}$ used in one instance where $\alpha_{k}$ should be).}

\textbf{Reply}: The notation has been corrected in the main text and appendix to account for this comment.

\subsection*{Comment 9}
{\it Figure 5: Because of the overlapping shading, I found this figure hard to read. One solution could be to use different line dashing/markers for the different overlaid density plots.}

\textbf{Reply}: Figure 5 has been amended and replaced by three plots side-by-side to improve readability.

\subsection*{Comment 10}
{\it In general, I think there are too many citations to the pesticide resistance, area-wide pest control, and invasive species literatures (related to my comment above that I think the comparison here is misplaced). }

\section*{Reviewer 2}
\subsection*{Comment 1}
{\it In the U.S. dairy industry, what is the approximate proportion of commercial dairy breeders who specialize in breeding and sell bull semen as the main business? And what is the proportion of dairy farmers who breed (most of) their cattle and raise them for dairy products? How does the split of the two types of businesses in the authors’ data set compared to that for the entire U.S. industry? Do the two types of business managers
differ significantly in their decisions on breeding (e.g., for the most profitable animal genetics to sell vs. for most productive animals to raise) and use of genomic testing (e.g., to grab market share vs. to avoid inbreeding)?}

\subsection*{Comment 2}
{\it Prior to the introduction of genomic testing, why were some lines of bulls more popular than others even though they had comparable or same genetic traits (page 15, top line)?

Was the popularity based upon prices, actual productivity (as opposed to the measured genetic traits), branding effects, or something else? After the introduction of genomic testing, did any already popular lines (e.g., super-sires) become less popular and vice versa (because buyers found out some already popular lines now have higher inbreeding rates)? Is it possible for breeders to line breed different lines of cattle since 2008 as some lines ascended to super-sires while other lines became less popular? In other words, can we guarantee there were no crossovers between the treatment group and the control group since 2008?}

\subsection*{Comment 3}
{\it Section 2 Conceptual framework examined the two key, sequential decisions facing animal breeders: to line breed or cross breed; to have genomic testing or not. It would be more informative to establish a unified model that explicitly addresses the interdependence between these two decisions with the objective function of market share maximization (say, to breed animals with the highest trait values relative to the market).
In this way, the authors can claim to explore the coordination between individual dairy breeders and the entire industry by analyzing the strategic actions of animal breeders.}

\subsection*{Comment 4}
{\it Regarding the cost-benefit analysis in section 6, the authors admitted that this is only a rough, back-of-the-envelop calculation, although there are still potentially substantial benefits of genomic testing to be included. Once animal breeders decide to have their animals genomically tested, they can ascertain the animals’ traits much sooner without
waiting an additional three or four years for the daughter-proven results. Meanwhile, animal breeders need to pay about \$50-100 per animal for testing and possibly cull the undesirable animals at a loss. The adjusted net benefits of genomic testing for the entire industry are likely to be greater than originally estimated, which may render inbreeding a lesser issue in terms of the economic value.}

\subsection*{Comment 5}
{\it I would also suggest adding more justifications when the authors try to project the estimated effects on inbreeding and net merit value using the authors’ sample data into the U.S. dairy industry. Related to my comment \#1 above, the adoption rates of genomic testing may vary between the (commercial) cattle breeders included in this study and cattle producers in the country, so the projected net benefits of genomic testing for the industry will have to be revised accordingly.}

\subsection*{Comment 6}
\begin{enumerate}
    \item Page 2, end of paragraph 1: “industr” should be “industry”.

    \textbf{Reply}: the typo was corrected.
    \item Page 7, line 4 from bottom: a period is missing before Melton et al.

    \textbf{Reply}: the typo was corrected.
    \item Page 12, 2nd paragraph, line 3 from the paragraph end: add some numbers to justify ``a substantial portion''.

    \textbf{Reply}: The paragraph now mentions that the NAAB represents 95\% of all livestock genetic companies in the United States.
    \item Page 14, 1st paragraph of section 4, lines 4-5: Why are bulls that are NOT genomically tested more likely to have better quality traits and be more widely used?

    \textbf{Reply}: The assertion is wrong; it was removed from the text.
    \item Page 16, line 4 from page bottom: Figure 5 in the main text is referred to here. Then drop ``in the Appendix''.

    \textbf{Reply}: the typo was corrected.
    \item Page 16, last sentence of the page: perhaps to conduct formal statistical tests on the similarities between the three distributions: matched, treat, and unmatched.

    \textbf{Reply}: A balance table was included, as a response to a previous comment.
    \item Page 24, line 3 below Figure 8: ``percent points'' should be ``percentage points''.
   
    \textbf{Reply}: the typo was corrected.
\end{enumerate}

\section*{Reviewer 3}
\subsection*{Comment 1}
{\it Does figure 2 have a source?  Also, what are units on inbreeding rate and how is it calculated? You use this inbreeding term quite a bit and it needs more explanation and definition.  Same for TPI on units and how calc in fig 2, though not as important for your paper as inbreeding rate. }

\subsection*{Comment 2}
{\it Page 6.  Your back of the envelope cost of inbreeding that you hang your hat on and note in the abstract of your paper is at much at \$6.7 billion. What is this cost? You mention Holstein USA but I do not see this reference in your reference list. When I review the publication at:  \url{https://www.holsteinusa.com/pdf/Interpreting_and_Utilizing_Haplotype_Information_1218.pdf} (which is in your references), I see no mention of the cost you calculate, but rather information about specific DNA sequences that are of concern. Then further, you state that the added milk productivity resulted in a net gain that exceeded this \$6.7 billion cost. So, net, genetic testing to date has been a huge benefit to the dairy industry if I am interpreting your claims correctly. So, why are you anchoring on this “cost” measure of \$6.7 billion.}

\subsection*{Comment 3}
{\it How much does genetic testing cost? Seems we need to know that to also assess net value.}

\textbf{Reply}: A single genetic test costs \$35. Source:
\url{https://vgl.ucdavis.edu/test/parentage-genetic-marker-report-cattle}.

\subsection*{Comment 4}
{\it Your premise that inbreeding is a problem because it “may have hidden costs” and not that it has been demonstrated to be an economic problem. That is, you rely on conjecture that inbreeding will result in major disease or other genetic concerns that lead I guess to important economic consequences because some traits are being ignored or are not exhibited sufficiently due to inbreeding. While this may be true from a genetics perspective, it seems like conjecture on your part at this point and has no economic assessment in your paper. The paper by Littledon (2024) \url{https://www.rroij.com/open-access/the-role-of-genomics-in-enhancing-dairy-cattle-productivity-and-health.pdf }  indicates genomics testing is enabling selecting for milk yield, feed efficiency, and detection of inherited deficiencies and mutations – essentially a more productive and healthy dairy herd. So, has inbreeding become problematic or do we know when it is a problem? What is the rate of inbreeding at which it is a problem? Are we at a critical level economically yet or are we a long ways away?}

\subsection*{Comment 5}
{\it Dairy producers are in business for the long run, so why would they act irrationally and select sires for only short-run gains and ignore disease resistance? Do you have any evidence that they actually do this? That is, breeders must provide genetics demanded by producers for long-run decisions. In the Bt corn case you discuss it was the seed companies that required refuge being planted so they could ensure the technology would be reliable in accomplishing what it claimed for the long run. This was not a federal government mandate. So, why would the same not occur in dairy genetics? You imply federal intervention is warranted because of an externality but you have not demonstrated such an externality exists or its magnitude. }

\subsection*{Comment 6}
{\it Do dairy producers not also conduct genetic tests increasingly on young heifers for deciding which ones to retain? If so, wouldn’t they also match that genetic info with sires for better complementarity of overall traits?}

\subsection*{Comment 7}
{\it Figure 4. Why has the number of bulls in the market gone up so rapidly from 2010 to 2015?  How does this mesh with your inbreeding measures? Are many of these bulls interrelated closely with each other?  At surface, it would seem more bulls would imply more genetic diversity in the market, but perhaps not?  Might deserve some explanation.}

\subsection*{Comment 8}
{\it Page 14 you state “Moreover, bulls that are not genomically tested are more likely to have better quality traits and be more widely used.”  Why would bulls that are not tested be more likely to have better quality traits?}

\textbf{Reply}: That assertion is wrong, and has been removed from the text.

\subsection*{Comment 9}
{\it Page 16, Figure 5 appears in the text, not the Appendix as stated here.}

\textbf{Reply}: The sentence was corrected to exclude references to the appendix.

\end{document}